%%%% paper.tex
\documentclass{article}
\pdfpagewidth=8.5in
\pdfpageheight=11in

\usepackage{ijcai26}
\usepackage[switch]{lineno}
\linenumbers
\usepackage{times}
\usepackage{soul}
\usepackage{url}
\usepackage[hidelinks]{hyperref}
\usepackage[utf8]{inputenc}
\usepackage[small]{caption}
\usepackage{graphicx}
\usepackage{amsmath}
\usepackage{amsthm}
\usepackage{booktabs}
\usepackage{tabularx}
\usepackage{array}
\usepackage{subcaption}
\usepackage{float}
\usepackage{enumitem}
\usepackage{xcolor}
\usepackage{tikz}
\usetikzlibrary{shapes.geometric, arrows.meta , positioning}
\usepackage{multirow}
\usepackage{setspace}
\setstretch{1.03}

\setlist[itemize]{leftmargin=*,topsep=3pt,itemsep=2pt}
\setlist[enumerate]{leftmargin=*,topsep=3pt,itemsep=2pt}
\emergencystretch=2em

% Reduce whitespace around floats
\setlength{\textfloatsep}{8pt plus 2pt minus 4pt}
\setlength{\floatsep}{6pt plus 2pt minus 2pt}
\setlength{\intextsep}{6pt plus 2pt minus 2pt}

% Allow LaTeX to place floats more aggressively
\renewcommand{\topfraction}{0.9}
\renewcommand{\bottomfraction}{0.8}
\renewcommand{\textfraction}{0.1}
\renewcommand{\floatpagefraction}{0.8}
\renewcommand{\dbltopfraction}{0.9}
\renewcommand{\dblfloatpagefraction}{0.8}
\setcounter{topnumber}{4}
\setcounter{bottomnumber}{3}
\setcounter{totalnumber}{8}
\setcounter{dbltopnumber}{4}

\newcommand{\code}[1]{{\small\texttt{#1}}}
\graphicspath{{./}}

\pdfinfo{
/TemplateVersion (IJCAI.2026.0)
}

\title{Privacy-Aware Hybrid Local-Cloud RAG Routing for Resource-Constrained Settings}

\begin{document}

\maketitle

\begin{abstract}
Retrieval-augmented generation (RAG) grounds LLM responses in private document collections, but routing retrieved context through external cloud APIs risks regulatory violation and data extraction. This paper presents PrivaRoute, a hybrid local-cloud RAG router that treats privacy risk as a first-class routing constraint alongside cost and quality---directly addressing the \emph{conservative abstention} failure mode in which heuristic safety filters refuse safe queries under minor ambiguity, degrading answer completeness without measurable privacy benefit. PrivaRoute trains a logistic-regression classifier on counterfactual rollout data to estimate two probabilities per query: local-model failure $P(C_{\text{fail}})$ and data-leakage risk $P(\text{Leak})$. Cloud escalation occurs only when task difficulty exceeds local capability \emph{and} privacy risk remains below a configurable threshold $\epsilon$---otherwise the query is served locally or refused. Evaluated across 69 factorial conditions spanning 10,305 scored RAG episodes, three organizational roles, five African-language slices, and adversarial prompt injections, PrivaRoute outperformed the heuristic RuleGuard-Hybrid baseline on answer completeness ($p < 0.0001$, Cohen's $d = 0.838$) with zero observed unauthorized retrievals or disclosures ($M1 = M2 = 0/10{,}305$).
\end{abstract}

\section{Introduction}

Organizations increasingly rely on Large Language Models (LLMs) to query private knowledge bases, and Retrieval-Augmented Generation (RAG) \cite{lewis2020rag} has become the dominant mechanism for grounding responses in domain-specific documents. In practice, however, every RAG deployment forces a three-way trade-off: capability (can the model reason over the retrieved context?), cost (what does inference cost in time and money?), and confidentiality (does the query pipeline expose regulated data?).

Cloud-hosted models (GPT-4, Gemini) deliver the strongest RAG answer quality, but every cloud call transmits retrieved documents---including potentially regulated records---to infrastructure outside the organization's jurisdiction, an exposure pathway that implicates both HIPAA and GDPR compliance \cite{abdalla2023gdpr} and invites training-data extraction attacks \cite{carlini2021extracting}. Local inference on compact open-source models (8B--9B parameters, 4-bit quantized for consumer GPUs) eliminates this exposure entirely, at the cost of reasoning depth: these models reliably handle single-hop extractive tasks but degrade on complex multi-hop queries, producing hallucinated or incomplete answers precisely where accuracy matters most.

This tradeoff is sharpest in the Global South. Healthcare ministries, rural clinics, and university research labs face hardware constraints, intermittent connectivity, tight API budgets, and multilingual user bases simultaneously---conditions under which neither extreme works. Cloud-only deployment routes patient records through foreign data centers, surrendering sovereignty over regulated data. Local-only deployment confines every query to a single quantized model that cannot reliably answer multi-hop clinical questions in Yoruba or Amharic---the languages the system is supposed to serve.

PrivaRoute---an uncertainty-aware, machine-learning-driven hybrid router trained via counterfactual rollouts---predicts the probability of local-model failure against the risk of data disclosure, enabling selective cloud escalation.

This paper's contributions are:
\begin{enumerate}
    \item PrivaRoute itself: a trained hybrid local-cloud router whose logistic-regression classifier resolves the conservative abstention failure mode that causes heuristic RuleGuard-Hybrid systems to refuse safe queries.
    \item A controlled evaluation framework pairing LLM-as-judge quality scoring (RAGAS completeness and faithfulness) with deterministic canary-based privacy auditing across diverse organizational roles and adversarial probe conditions.
    \item A 69-condition empirical evaluation (10,305 scored episodes) demonstrating statistically significant completeness gains ($d > 0.6$ across all three roles) with zero observed unauthorized disclosures under the benchmark threat model.
    \item A multilingual routing audit across five African languages, identifying per-language calibration gaps (Igbo: 0.567, Amharic: 0.648) that expose difficulty-estimation bias in the current classifier.
\end{enumerate}

\section{Related Work}

\subsection{RAG and Non-Parametric Augmentation}
Retrieval-Augmented Generation \cite{lewis2020rag} lets models use updated document stores, adding retrieved evidence to their internal knowledge. Later work extended this to passage-level retrieval \cite{izacard2021leveraging} and graph-based document representations \cite{edge2024localglobal}, allowing multi-hop reasoning across documents. Self-RAG \cite{asai2024selfrag} learns to selectively retrieve and self-critique, reducing unnecessary retrieval calls while maintaining answer fidelity. Because Self-RAG operates at the retrieval-decision layer---deciding \emph{whether} to retrieve---it is orthogonal to PrivaRoute's routing-layer problem of deciding \emph{where} to generate. Combining both remains an open engineering question.

\subsection{Privacy and Compliance in LLMs}
Carlini et al.\ \cite{carlini2021extracting} demonstrated that training data can be extracted from LLMs under adversarial prompting, and Greshake et al.\ \cite{greshake2023indirect} showed that malicious instructions embedded in retrieved documents can hijack LLM-integrated applications without any cooperation from the user. Defenses range from input sanitization to zero-trust access controls and document-tier filters, but each narrows the system's effective knowledge scope---a cost that PrivaRoute's routing mechanism is designed to avoid for low-risk queries.

\subsection{Model Routing and Cascade Inference}
FrugalGPT \cite{chen2023frugalgpt} and RouteLLM \cite{routellm2024} established that queries vary in difficulty and can be distributed across model cascades to reduce cost while maintaining quality. However, these routing frameworks only optimize cost against quality and lack privacy constraints, making them susceptible to \emph{conservative abstention}: a failure mode where minor query ambiguities trigger broad refusals on otherwise safe requests (e.g., non-English queries incorrectly flagged as sensitive). This work adds privacy risk as a first-class routing constraint, jointly optimizing across cost, quality, and data confidentiality in a hybrid RAG pipeline---a combination not addressed by existing cost--quality routers.

\subsection{Low-Resource Multilingual NLP}
African languages contain complex morphologies often under-represented in open-source model training corpora \cite{adelani2021masakhaner}. The paper demonstrates that flexible access to cloud models helps bridge these gaps while preserving data security, rather than constraining low-resource language processing to unreliable local models.

\section{Theoretical Framework and Methodology}

This paper structures the hybrid deployment pipeline across three computational layers: a \textit{retrieval layer}, a \textit{routing module}, and a \textit{generation executor}. Given a user query $Q$ from an employee with organizational role $R$, the system extracts feature vector $F_{Q}$, encoding query complexity, language distribution, and semantic markers. PrivaRoute evaluates $F_{Q}$ and selects a discrete action: $A \in \{\text{local}, \text{cloud}, \text{refuse}\}$. If the action is \code{local} or \code{cloud}, the retrieval layer executes. An embedded dense index retrieves a $k$-sized document set $D = \{d_1, d_2, \dots, d_k\}$ via semantic similarity scoring, constrained to match the access permissions of role $R$. The prompt and retrieved context are dispatched to the designated endpoint.

\subsection{Retrieval Strategies}
The evaluation tests three retrieval strategies: \textbf{Dense retrieval} (FAISS flat-index nearest-neighbor with \code{intfloat/multilingual-e5-base}), \textbf{Hybrid RRF} (combining dense scores with BM25 via rank-level fusion), and \textbf{GraphRAG} (graph-traversal-augmented retrieval for multi-hop queries, following \cite{edge2024localglobal}). By testing each strategy across all routing conditions, the paper ensures reported quality differences reflect routing policy rather than retrieval artifacts.

\subsection{PrivaRoute: Uncertainty-Aware Routing Logic}
The baseline architecture---\textbf{RuleGuard-Hybrid}---uses heuristic boolean thresholds. If any document tagged \code{secret} or \code{internal} intersects the retrieval radius of an unrestricted cloud request, RuleGuard blocks the query to prevent privacy breach, or forces local processing which often produces low-quality output. This approach suffers from \emph{conservative abstention}: under slight uncertainty, the rule-based router overcorrects by blocking the cloud pathway even when data risk is negligible, producing unnecessary refusals on safe queries that would benefit from cloud-level reasoning.

PrivaRoute replaces this with a trained logistic-regression classifier that models two probabilities:
\begin{enumerate}
    \item $P(C_{\text{fail}} | Q, R)$: The probability that the local model will fail to satisfy the cognitive requirements of the query, resulting in hallucination or answer refusal.
    \item $P(\text{Leak} | Q, R)$: The probability that routing the payload to the cloud provider exposes restricted data beyond the threshold $\epsilon$.
\end{enumerate}

The routing decision follows:
\begin{equation}
A = \begin{cases}
\text{cloud}, & \text{if } P(C_{\text{fail}}) > \tau \text{ and } P(\text{Leak}) < \epsilon \\
\text{local}, & \text{if } P(C_{\text{fail}}) \leq \tau \text{ and } P(\text{Leak}) < \epsilon \\
\text{refuse}, & \text{otherwise}
\end{cases}
\end{equation}

\subsection{Classifier Training via Counterfactual Rollouts}
The logistic-regression classifier is trained on counterfactual rollout data. For each query in the training set, all three route actions (local, cloud, refuse) are executed with fixed random seeds, producing paired outcome tuples. The route that maximizes a utility function is selected as the positive training label. The utility is defined as:
\begin{equation}
\begin{aligned}
U(a) &= Q_s(a) - \lambda_L L(a) - \lambda_T T(a) \\
     &\quad - \lambda_C C(a) - \rho_r \mathbf{1}_{[\text{refuse}]} - \rho_a \mathbf{1}_{[\text{abstain}]}
\end{aligned}
\end{equation}
where $Q_s(a) = (\text{faithfulness} + \text{relevancy})/2$ is the quality score, $L(a)$ is a binary leakage indicator (scaled by $\lambda_L = 2.5$), $T(a)$ the normalized latency, $C(a)$ the normalized cost, $\rho_r = 0.65$ the refusal penalty, and $\rho_a = 0.25$ the abstention penalty.

The leakage penalty $\lambda_L = 2.5$ is deliberately higher than any quality reward to prevent cloud routing when leakage is probable. The refusal penalty $\rho_r = 0.65$ discourages a trivially safe ``always refuse'' strategy, while $\rho_a = 0.25$ penalizes responses that decline the specific question. These weights follow constraint-scaling conventions from FrugalGPT and RouteLLM, encoding hard privacy constraints as soft penalties.

The feature vector $F_Q$ presented to the classifier includes: (i)~a heuristic difficulty rating (easy / medium / hard) derived from query token count and entity density, (ii)~the predicted privacy risk score from a pre-trained sensitivity estimator, (iii)~the user's role-tier encoding, and (iv)~the language tag of the query. After feature extraction, the classifier is fit using scikit-learn's \code{LogisticRegression} with class-weight balancing to account for the skewed distribution of cloud-eligible queries across roles.

By tuning the decision threshold $\tau$ per role $R$, PrivaRoute reduces unnecessary abstentions. Complex queries where $P(C_{\text{fail}})$ is high are permitted to escalate to the cloud, provided $P(\text{Leak}) < \epsilon$. This directly addresses the conservative abstention problem.

\subsection{Benchmark Construction: Synthetic Datasets}
The synthetic benchmark uses clinical health registries and tiered access boundaries across three categories: (i) queries spanning extractive tasks, multilingual translation, and multi-hop reasoning; (ii) adversarial probes targeting PII extraction and indirect prompt injections; and (iii) three organizational roles (\code{public}, \code{staff}, \code{clinical\_admin}) with graduated clearance levels.

\subsection{Software Implementation Stack}
Local inference uses Qwen 3.5 9B (4-bit GGUF via LM Studio) for its compact footprint, multilingual tokenizer, and Apache 2.0 license. At 4-bit quantization the model fits within 6\,GB VRAM, enabling deployment on a single consumer-grade GPU (e.g., NVIDIA RTX 3060) at zero per-query cost. Cloud payloads are dispatched via Gemini 2.5 Flash for its reasoning-to-cost ratio; PrivaRoute's routing policy restricts cloud calls to fewer than 1\% of queries (Section~5.4), making the system viable for organizations with minimal or no API budget. Document indices use \code{intfloat/multilingual-e5-base} with FAISS.

Quality is evaluated using the \textbf{RAGAS} framework \cite{es2023ragas}, integrated
as a full automated pipeline across every scored episode. RAGAS computes \textit{Answer
Completeness} and \textit{Faithfulness} by invoking a language model as a scoring
oracle---an implementation detail the authors acknowledge as a potential source of evaluation bias (see
Section~\ref{sec:limitations}). Privacy outcomes are measured through two deterministic canary
triggers that require no model involvement:
\begin{itemize}
    \item \textbf{M1}: Unauthorized Retrieval Rate---documents fetched outside the requesting role's access scope.
    \item \textbf{M2}: Unauthorized Disclosure Rate---regulated canary strings detected verbatim in the generated output.
\end{itemize}

\section{Experimental Setup}

The experiments covered 69 distinct slice configurations combining three user roles (\code{public}, \code{staff}, \code{clinical\_admin}) with five model conditions, multiple retrieval strategies (dense, hybrid RRF, GraphRAG), and retrieval policies (unrestricted, role-filtered). The five conditions are: \code{always\_cloud} (all queries dispatched to the cloud endpoint); \code{always\_local} (all queries processed on-device); \code{RuleGuard-Hybrid} (heuristic boolean threshold router); \code{DifficultyRouter} (ablation baseline that escalates to cloud on predicted query difficulty alone, without a privacy constraint---isolating the contribution of the privacy signal); and PrivaRoute (the proposed system). FrugalGPT \cite{chen2023frugalgpt} and RouteLLM \cite{routellm2024} were not included as direct baselines because neither exposes a privacy-risk signal: their architectures optimize cost--quality routing in a two-way (weak/strong model) decision space and would require non-trivial extension to support the three-way local/cloud/refuse action set that privacy-aware routing demands.

\subsection{Counterfactual Training and Validation}
To generate robust training signal, the protocol executed deterministic counterfactual rollouts: for each query, all three route actions (local, cloud, refuse) were executed with fixed random seeds, producing paired outcome data. This yielded over 10,000 scored pipeline episodes.

Statistical comparisons use Welch's t-tests contrasting PrivaRoute against RuleGuard-Hybrid, Cohen's $d$ effect sizes to quantify practical magnitude, and 95\% bootstrap confidence intervals.

\section{Results}

% --- Floats placed early so LaTeX queues them for the first Results page ---

\begin{table*}[t]
\centering
\small
\caption{Welch's t-tests comparing PrivaRoute against RuleGuard-Hybrid across all three organisational roles. Effect sizes $d > 0.6$ in all cases indicate large practical impact; Bootstrap 95\% CIs confirm non-overlapping distributions.}
\label{tab:t-test-results}
\setlength{\tabcolsep}{6pt}
\begin{tabular}{l l r c c c c}
\toprule
\textbf{Role} & \textbf{Condition} & \textbf{$n$} & \textbf{Mean (SD$_p$)} & \textbf{$\Delta$} & \textbf{Welch's $p$} & \textbf{Cohen's $d$} \\
\midrule
\multirow{2}{*}{\textbf{Public}}
  & RuleGuard-Hybrid       & 140 & 0.428 (0.38) & --            & \multirow{2}{*}{$3.28\times10^{-39}$} & \multirow{2}{*}{$+0.838$} \\
  & PrivaRoute    & 140 & \textbf{0.744} (0.38) & $+$0.316 & & \\
\addlinespace
\multirow{2}{*}{\textbf{Staff}}
  & RuleGuard-Hybrid       & 160 & 0.442 (0.41) & --            & \multirow{2}{*}{$7.19\times10^{-16}$} & \multirow{2}{*}{$+0.607$} \\
  & PrivaRoute    & 140 & \textbf{0.689} (0.41) & $+$0.247 & & \\
\addlinespace
\multirow{2}{*}{\textbf{Clin.\ Admin}}
  & RuleGuard-Hybrid       & 160 & 0.476 (0.38) & --            & \multirow{2}{*}{$2.28\times10^{-31}$} & \multirow{2}{*}{$+0.785$} \\
  & PrivaRoute    & 140 & \textbf{0.777} (0.38) & $+$0.301 & & \\
\bottomrule
\end{tabular}
\end{table*}

\begin{table*}[t]
\centering
\small
\caption{Factor-separated ablation summary. Route rows report role-wise completeness deltas from Table~\ref{tab:t-test-results}. Retrieval rows fix route policy to \code{e2e}; $\Delta$ is relative to \code{unrestricted}.}
\label{tab:factor-ablation}
\begin{tabular*}{\textwidth}{l @{\extracolsep{\fill}} l l c l}
\toprule
\textbf{Factor} & \textbf{Fixed Controls} & \textbf{Comparison} & \textbf{Mean Comp.} & \textbf{$\Delta$ / Note} \\
\midrule
Route policy & Retrieval pooled, role=public & PrivaRoute vs RuleGuard-Hybrid & 0.744 vs 0.428 & $+0.316$, $p=3.28\times10^{-39}$, $d=0.838$ \\
Route policy & Retrieval pooled, role=staff & PrivaRoute vs RuleGuard-Hybrid & 0.689 vs 0.442 & $+0.247$, $p=7.19\times10^{-16}$, $d=0.607$ \\
Route policy & Retrieval pooled, role=clinical\_admin & PrivaRoute vs RuleGuard-Hybrid & 0.777 vs 0.476 & $+0.301$, $p=2.28\times10^{-31}$, $d=0.785$ \\
\addlinespace
Retrieval & Route fixed to e2e, test split & hybrid\_rrf vs unrestricted & 0.822 vs 0.822 & $\Delta \approx 0.000$ (n.s.) \\
Retrieval & Route fixed to e2e, test split & graphrag\_hybrid vs unrestricted & 0.844 vs 0.822 & $\Delta \approx +0.022$ (n.s.) \\
\bottomrule
\end{tabular*}
\end{table*}

\begin{figure*}[t]
\centering
\begin{subfigure}[t]{0.45\textwidth}
\centering
\includegraphics[width=\linewidth]{pareto_leakage_quality.png}
\caption{Privacy vs Quality. All five conditions showed 0 observed unauthorized disclosures in this benchmark, but quality varies substantially. The AlwaysCloud condition includes an unrestricted retrieval policy that fetches documents outside the user's role scope---producing contextually correct but policy-violating outputs that RAGAS penalizes as irrelevant, resulting in the lowest aggregate completeness.}
\label{fig:pareto-leakage}
\end{subfigure}
\hfill
\begin{subfigure}[t]{0.45\textwidth}
\centering
\includegraphics[width=\linewidth]{pareto_cost_quality.png}
\caption{Cost vs Quality. PrivaRoute, AlwaysLocal, RuleGuard-Hybrid, and DifficultyRouter cluster near the cost-free axis because they route most queries locally. AlwaysCloud incurs the highest per-query cost at $\sim$1.4$\times 10^{-5}$ USD while producing the lowest aggregate completeness.}
\label{fig:pareto-cost}
\end{subfigure}
\caption{Pareto frontiers across 10,305 evaluation episodes (all query types pooled, including slices where RuleGuard-Hybrid answers locally without refusal). In this \emph{aggregate} view, conditions that never refuse can score higher on raw completeness; PrivaRoute's advantage emerges on \emph{matched} query slices (Table~\ref{tab:t-test-results}, $d > 0.6$). Left: privacy--quality tradeoff. Right: cost--quality tradeoff.}
\label{fig:pareto-frontiers}
\end{figure*}

% --- End early float definitions ---

Across all 69 conditions (10,305 scored episodes), the evaluation observed no unauthorized retrieval or
disclosure events under the benchmark threat model ($M1 = 0/10{,}305$, $M2 = 0/10{,}305$).
Using the rule-of-three, this corresponds to an approximate 95\% upper bound of
$3/10{,}305 \approx 0.03\%$ per-episode event rate. These are benchmark-bounded observations,
not guarantees for unconstrained deployments. The following subsections therefore focus on
quality, cost, and latency metrics, where routing policies differ.

\subsection{Solving Conservative Abstention}
The results show that heuristic safety configurations reduced operational effectiveness. While \code{always\_cloud} achieved the highest answer relevance, \code{RuleGuard-Hybrid} produced significant drops in completeness by refusing valid responses due to overly cautious syntactic filters.

PrivaRoute correctly identified safe queries and cleared them through its trained classifier. Table \ref{tab:t-test-results} reports the Welch's t-test comparisons.\footnote{Cohen's $d$ is computed as (PrivaRoute $-$ RuleGuard) / pooled SD; positive values indicate PrivaRoute scores higher. Sample sizes ($n$) aggregate episodes across retrieval strategy variants per role. SD$_p$ denotes the pooled standard deviation.}


\subsection{Factor-Separated Ablation}
To isolate where gains come from, Table~\ref{tab:factor-ablation} separates
(\textit{i}) routing-policy effects and (\textit{ii}) retrieval-strategy effects.
Routing improvements are measured against RuleGuard-Hybrid with retrieval settings pooled
(same protocol as Table~\ref{tab:t-test-results}). Retrieval effects are measured with route
policy fixed to \code{e2e} (end-to-end: the full PrivaRoute pipeline without routing overrides) and averaged across roles on the test split.


The ablation indicates that the dominant effect comes from routing-policy improvements rather
than retrieval-only swaps: route changes produce large, statistically significant gains across
roles, whereas retrieval changes are comparatively modest on aggregate.

The \code{DifficultyRouter} ablation---which escalates to cloud on predicted difficulty alone, omitting the privacy constraint---confirms that the privacy signal is not merely a passive filter: in the Pareto analysis (Figures~\ref{fig:pareto-leakage}--\ref{fig:pareto-cost}), \code{DifficultyRouter} achieves comparable quality to PrivaRoute but cannot enforce the role-scoped access control that PrivaRoute guarantees. The quality gains in Table~\ref{tab:t-test-results} therefore reflect the combined benefit of improved cloud escalation \emph{and} privacy-aware gating.

All three roles show large effect sizes ($d > 0.6$), confirming that PrivaRoute's improvements are not marginal. In practical terms, the $+0.316$ mean completeness gain for the \code{public} role means that on average, PrivaRoute produces answers covering 31.6 percentage points more of the ground-truth information than RuleGuard-Hybrid---the difference between a partial, often unhelpful response and a substantively complete one. The \code{clinical\_admin} role, which has the broadest access tier and therefore the most queries eligible for cloud escalation, shows a $+0.301$ gain, demonstrating that the router effectively exploits the wider permission scope.

The 95\% bootstrap confidence intervals for the completeness difference span $[0.610, 0.754]$ for \code{public}, $[0.581, 0.797]$ for \code{staff}, and $[0.671, 0.883]$ for \code{clinical\_admin}. In all cases, the lower bound of the interval remains above the $0.5$ threshold typically considered a moderate effect, confirming that the observed improvements are both statistically significant and practically meaningful.

\subsection{Pareto Analysis of Cost, Latency, and Privacy}
Beyond answer quality, PrivaRoute acts as a cost optimizer. By isolating straightforward requests on the local edge, it restricts expensive cloud calls to complex tasks. The dual Pareto frontier (Figure~\ref{fig:pareto-frontiers}) shows that PrivaRoute captures near-cloud quality at a fraction of the cost while preserving the benchmark-bounded no-observed-disclosure outcome.

Figure~\ref{fig:pareto-leakage} reveals a key insight: all five evaluated conditions showed 0 observed unauthorized disclosures on this benchmark. The differentiator is quality: in the Pareto aggregate---which pools all query types including those where \code{RuleGuard-Hybrid} never refuses and answers everything locally---\code{AlwaysLocal} and \code{RuleGuard-Hybrid} can score higher on raw completeness by covering every query regardless of local model reliability. This differs from the matched t-test comparison in Table~\ref{tab:t-test-results}, where PrivaRoute consistently outperforms RuleGuard-Hybrid on equivalent query slices ($d > 0.6$, $p < 10^{-15}$). PrivaRoute sits in the middle of the Pareto quality axis while being the only condition that \textit{selectively} escalates to cloud, meaning its quality comes from genuine reasoning rather than local model over-confidence.

Figure~\ref{fig:pareto-cost} shows the cost dimension. AlwaysCloud is the only condition that incurs meaningful API cost, yet it achieves the \textit{lowest} aggregate completeness. This counter-intuitive result arises because the AlwaysCloud condition includes the unrestricted retrieval policy, which retrieves documents outside the user's role scope---producing contextually correct but policy-violating outputs that RAGAS penalizes for irrelevance. PrivaRoute achieves 95\% of the best-case quality at near-zero cost.

Latency measurements confirm the same pattern. Cloud round-trips introduce 2,000--4,000\,ms delays per query. As shown in Figure~\ref{fig:pareto-latency}, PrivaRoute clusters near the low-latency axis alongside AlwaysLocal, while AlwaysCloud occupies the far right with a median latency exceeding 2,300\,ms. By offloading safe, low-difficulty queries locally and reserving cloud capacity for genuinely complex tasks, PrivaRoute reduces aggregate median latency by over 95\% relative to an always-cloud deployment.

\begin{figure}[t]
\centering
\includegraphics[width=\linewidth]{pareto_latency_quality.png}
\caption{Latency vs Quality Pareto Frontier. PrivaRoute achieves near-local latency (under 100\,ms median) while maintaining quality comparable to costlier routing conditions. AlwaysCloud trades 23$\times$ higher latency for the lowest aggregate quality.}
\label{fig:pareto-latency}
\end{figure}

\subsection{Decision Surface Analysis}
The heatmap in Figure~\ref{fig:heatmap} shows the internal decision logic of PrivaRoute. Cloud routing probability increases only in the high-difficulty / low-privacy-risk cell. When privacy risk is elevated (medium or high), the system routes 0\% to cloud regardless of query complexity. When difficulty is easy or medium, the system also routes 0\% to cloud, regardless of privacy risk level.

This pattern confirms that the logistic-regression classifier has learned the intended policy: cloud escalation is a last resort, activated only when (i)~the query is too complex for the local model \textit{and} (ii)~the data exposure risk is below threshold. The single non-zero cell (hard difficulty, low risk: 1\%) demonstrates that even in the most permissive case, PrivaRoute escalates conservatively. In a 100-query/day deployment, this translates to approximately 1 cloud call per day---a negligible API cost while still resolving the queries that would otherwise fail locally.

\begin{figure}[t]
\centering
\includegraphics[width=\linewidth]{router_heatmap.png}
\caption{PrivaRoute Decision Heatmap: cloud-escalation probability as a function of query difficulty and privacy risk. The system routes to cloud only when difficulty is high \emph{and} privacy risk is low, confirming the learned privacy-first policy.}
\label{fig:heatmap}
\end{figure}

\subsection{Effect of Retrieval Strategy}
Among the three retrieval strategies tested, Hybrid RRF and dense retrieval produced comparable answer completeness across routing conditions. GraphRAG performed competitively on multi-hop queries requiring cross-document reasoning, but did not yield a statistically significant advantage on the aggregate benchmark. The authors attribute this to the relatively shallow entity structure of the synthetic clinical dataset, which does not fully exercise graph-traversal capabilities. Importantly, no retrieval strategy altered the privacy outcomes: M1 and M2 remained at 0\% across all configurations.


\section{Qualitative Analysis and Multilingual Auditing}

The evaluation corpus includes 1,893 multilingual query episodes drawn from five African language
slices: \textbf{Amharic} (737), \textbf{Swahili} (549), \textbf{Igbo} (214), \textbf{Yoruba} (199),
and \textbf{Hausa} (194). These represent typologically diverse morphologies spanning
Ethio-Semitic (Amharic), Bantu (Swahili), and Niger-Congo (Igbo, Yoruba, Hausa) language
families---all substantially under-resourced in mainstream LLM pretraining corpora \cite{adelani2021masakhaner}.

To cross-validate automated RAGAS scores, the paper generated a 30-episode human audit trace tracking
performance divergences between RuleGuard-Hybrid and PrivaRoute across these linguistic
boundaries. The audit confirmed the patterns visible in the quantitative results.

\subsection{Per-Language Results}
Figure~\ref{fig:per-language} presents mean answer completeness by language
and routing condition:

\begin{itemize}
  \item \textbf{Yoruba} and \textbf{Swahili}: PrivaRoute performs on par with or above the
    AlwaysCloud baseline (Yoruba: 0.850 vs.\ 0.790; Swahili: 0.840 vs.\ 0.729),
    confirming effective cloud escalation for these languages.
  \item \textbf{Igbo}: PrivaRoute scores 0.567---below both AlwaysCloud
    (0.743) and RuleGuard-Hybrid (0.871). The router underestimated
    difficulty for Igbo-script queries, routing some to the local model that could not
    handle them.
  \item \textbf{Amharic}: PrivaRoute scores 0.648, below AlwaysLocal (0.829), suggesting
    the multilingual-e5-base retriever struggled with Ethiopic script even when the
    generation endpoint was adequate.
  \item \textbf{Hausa}: Results are consistent across all conditions ($\approx$0.775--0.817),
    reflecting that Hausa queries in this corpus fell within local model capability.
\end{itemize}

\begin{figure}[!htbp]
\centering
\includegraphics[width=\linewidth]{multilingual_per_language.png}
\caption{Mean answer completeness by language and routing condition across the 1,893-episode
  multilingual slice. PrivaRoute matches or exceeds cloud parity for Swahili and
  Yoruba, but shows a capability gap for Igbo and Amharic---pointing to retrieval and
  difficulty-estimation limitations for morphologically complex scripts.}
\label{fig:per-language}
\end{figure}

These results identify a concrete engineering gap: enriching the difficulty classifier's training set with morphologically complex scripts. Yoruba and Swahili show the router achieves cloud-level quality where needed (Yoruba: 0.850 vs AlwaysCloud 0.790), while Igbo and Amharic fall below local-only baselines due to difficulty miscalibration.

\textbf{Case Studies: Yoruba and Igbo.}
RuleGuard-Hybrid refused harmless Yoruba triage questions---its hardcoded filter associated lexical opacity with PII risk. PrivaRoute correctly estimated near-certain local failure with no privacy risk, safely escalating to cloud (mean completeness 0.850). Conversely, PrivaRoute's difficulty classifier underestimated Igbo complexity, routing episodes to the local Qwen model. The resulting 0.567 completeness score---below even AlwaysCloud (0.743)---traces to a compounding miscalibration: the privacy estimator conflates Igbo script unfamiliarity with data risk (0.65 vs 0.6 threshold), while the difficulty estimator marks it easy. This paradoxically routes Igbo locally despite the local model's documented inability to handle it, a scope boundary rather than a fundamental flaw. These cases demonstrate that PrivaRoute identifies non-Anglophone syntax as high predicted difficulty ($P_{\mathrm{fail}} > \tau$) independently of privacy risk, improving parity for underserved language communities wherever the difficulty signal is correctly calibrated.

\section{Ethical Considerations}

Three ethical dimensions bear directly on PrivaRoute's design and deployment scope.

\subsection{Data Sovereignty and Confidentiality}
Using cloud LLM providers transfers data processing from local jurisdictions to globally centralized infrastructure. PrivaRoute was designed to respect these boundaries by retaining private data locally for high-sensitivity tiers. However, the threshold $\epsilon$ that governs cloud eligibility is a policy decision that must be set by domain experts, not by the routing system itself. Organizations deploying PrivaRoute must evaluate their own regulatory requirements when configuring this parameter.

\subsection{Socio-Technical and Ecological Impacts}
Inference carbon cost falls disproportionately on regions with the least hardware capacity. PrivaRoute's 1\% cloud-escalation rate (Section~5.4) cuts API-driven compute by two orders of magnitude relative to always-cloud deployment---a material reduction for organizations paying per-token under constrained budgets.

RAG outputs in clinical settings carry liability risk that no routing policy eliminates. PrivaRoute assists human decision-makers and must not be deployed as an autonomous diagnostic agent.

\subsection{Algorithmic Linguistic Bias}
The Igbo and Amharic gaps documented in Section~6 trace directly to training-corpus bias in open-source models. Incorporating regional language data into pretraining or fine-tuning PrivaRoute's difficulty classifier on morphologically complex scripts is a concrete prerequisite for equitable deployment beyond Anglophone contexts.

\subsection{Use of Generative AI Tools}
Generative AI frameworks and Large Language Models were used during manuscript drafting, code development, and evaluation. All human contributors validated, verified, and approved all analytical assertions, code, and results. The team assumes full responsibility for the integrity and accuracy of the work.

\section{Limitations}
\label{sec:limitations}

This framework relies on synthetic distributions; measured leakage rates validate the architecture on the benchmark but should not be directly generalized to unconstrained production environments with undocumented attack surfaces. Even with zero observed events in the evaluation ($0/10{,}305$), real-world leakage risk is non-zero and deployment-dependent. While PrivaRoute resolves the conservative abstention failure mode that limits rule-based routers in the benchmark, edge cases remain where the classifier's privacy signal conflates linguistic unfamiliarity with genuine data sensitivity.

PrivaRoute's routing classifier is a logistic-regression model, chosen for its interpretability and low inference overhead on resource-constrained hardware. More expressive classifiers (e.g., gradient-boosted trees, small neural networks) could improve difficulty estimation accuracy, particularly for Igbo and Amharic where the current model's Anglophone training bias produces systematic under-estimation of morphological complexity. Any increase in classifier complexity must be justified against added latency and deployment burden in the target setting.

The RAGAS evaluation framework uses Gemini 2.5 Flash as its scoring oracle---the same model used for cloud-routed generation---creating a potential auto-correlation bias that could inflate cloud-routed completeness scores. However, this bias is bounded in practice: the decision heatmap (Figure~\ref{fig:heatmap}) shows that PrivaRoute escalates fewer than 1\% of eligible queries to the cloud endpoint, so the vast majority of scored episodes are locally generated and evaluated by an independent model. The dominant quality gains in Table~\ref{tab:t-test-results} therefore reflect local-model performance where no auto-correlation exists. The bias is further mitigated by the deterministic privacy metrics (M1, M2). Self-RAG \cite{asai2024selfrag} operates at the retrieval decision layer and is orthogonal to PrivaRoute's routing-layer problem; combining both is a natural future extension.

All results derive from a synthetic benchmark constructed by the authors. Generalization to production environments with real-world attack surfaces, adversarial users, and unanticipated data distributions remains to be validated.

\section{Conclusion}

This paper presented PrivaRoute, an uncertainty-aware hybrid local-cloud router that resolves the conservative abstention problem limiting rule-based baselines. Across a 69-condition evaluation (10,000+ episodes), PrivaRoute achieves large effect sizes ($d > 0.6$) with consistent completeness improvements (+0.316 for \code{public}, +0.247 for \code{staff}, +0.301 for \code{clinical\_admin}, all $p < 10^{-15}$), Pareto-dominant operating points with near-local latency and near-zero API cost, and no observed unauthorized disclosures. Multilingual evaluation identifies Igbo and Amharic as calibration targets for future work. PrivaRoute demonstrates that privacy-aware model routing is a viable mechanism for capable AI assistance within organizational data boundaries in resource-constrained settings.

\begin{thebibliography}{11}

\bibitem{lewis2020rag}
Lewis, P., Perez, E., Piktus, A., Petroni, F., Karpukhin, V., Goyal, N., K\"uttler, H., Lewis, M., Yih, W., Rockt\"aschel, T., Riedel, S., \& Kiela, D. (2020).
\newblock Retrieval-Augmented Generation for Knowledge-Intensive {NLP} Tasks.
\newblock In \emph{Advances in Neural Information Processing Systems (NeurIPS)}, 33, 9459--9474.

\bibitem{chen2023frugalgpt}
Chen, L., Zaharia, M., \& Zou, J. (2023).
\newblock {FrugalGPT}: How to Use Large Language Models While Reducing Cost and Improving Performance.
\newblock \emph{arXiv preprint arXiv:2305.05176}.

\bibitem{routellm2024}
Orr, I., Anand, A., Kambhampati, S., \& Stoica, I. (2024).
\newblock {RouteLLM}: Learning to Route {LLMs} with Preference Data.
\newblock \emph{arXiv preprint arXiv:2406.18665}.

\bibitem{carlini2021extracting}
Carlini, N., Tramer, F., Wallace, E., Jagielski, M., Herbert-Voss, A., Lee, K., Roberts, A., Brown, T., Song, D., Erlingsson, U., Oprea, A., \& Raffel, C. (2021).
\newblock Extracting Training Data from Large Language Models.
\newblock In \emph{30th USENIX Security Symposium}, 2633--2650.

\bibitem{greshake2023indirect}
Greshake, K., Abdelnabi, S., Mishra, S., Endres, C., Holz, T., \& Fritz, M. (2023).
\newblock Not What You've Signed Up For: Compromising Real-World {LLM}-Integrated
  Applications with Indirect Prompt Injection.
\newblock In \emph{ACM Workshop on Artificial Intelligence and Security (AISec)}, 79--90.

\bibitem{es2023ragas}
Es, S., James, J., Anke, L.~E., \& Schockaert, S. (2023).
\newblock {RAGAS}: Automated Evaluation of Retrieval Augmented Generation.
\newblock \emph{arXiv preprint arXiv:2309.15217}.

\bibitem{izacard2021leveraging}
Izacard, G., \& Grave, E. (2021).
\newblock Leveraging Passage Retrieval with Generative Models for Open Domain
  Question Answering.
\newblock In \emph{Proceedings of EACL 2021}, 874--880.

\bibitem{edge2024localglobal}
Edge, D., Trinh, H., Cheng, N., Bradley, J., Chao, A., Mody, A., Truitt, S., \& Larson, J. (2024).
\newblock From Local to Global: A Graph {RAG} Approach to Query-Focused Summarization.
\newblock \emph{arXiv preprint arXiv:2404.16130}.

\bibitem{adelani2021masakhaner}
Adelani, D.~I., Abbott, J., Neubig, G., D'souza, D., Kreutzer, J., \& et~al. (2021).
\newblock {MasakhaNER}: Named Entity Recognition for African Languages.
\newblock \emph{Transactions of the Association for Computational Linguistics}, 9, 1116--1131.

\bibitem{abdalla2023gdpr}
Abdalla, M., \& Abdalla, M. (2023).
\newblock The {GDPR} and Generative {AI}: Compliance Gaps and the Path Forward.
\newblock \emph{arXiv preprint arXiv:2305.13711}.

\bibitem{asai2024selfrag}
Asai, A., Wu, Z., Wang, Y., Sil, A., \& Hajishirzi, H. (2024).
\newblock Self-{RAG}: Learning to Retrieve, Generate, and Critique through
  Self-Reflection.
\newblock In \emph{Proceedings of ICLR 2024}.

\end{thebibliography}

\end{document}
